\chapter{Summary and future directions} \label{ch:summary}

The goals of this book as outlined in the introductory chapter are to improve the state of description of the grammars of Chadic languages, contribute to the typology of Chadic and Afroasiatic languages, and to provide insights into the nature of how motion is expressed through linguistic structures.

\section{Chadic languages}

The concepts outlined in this book, especially the types of directional meaning (Chapter \ref{ch:distribution}) and the typology of co-expression with directional meaning (Chapter \ref{ch:functions}) provide a framework for describing and comparing this frequently attested element of Chadic verbal morphology. The use of consistent comparative concepts applied to each of the Chadic languages covered in this book have made their descriptions more cohesive and accessible, and the process also provided an occasion to review and correct some of the inconsistencies found in past publications (as detailed in the appendices). 

One improvement needed in future descriptive work is a more explicit and consistent analyses of the function of directional extensions with all kinds of verbs. Our understanding of how speakers of Chadic languages use directional extensions is significantly limited due to the fact that in half of the examples of directional extensions it is not clear what the function is (Chapter \ref{ch:functions}). It is possible that many of the translations that lack an explicit translation of the directional extension's function are cases of aspectual uses of directional extensions. This is a topic that deserves more attention in future descriptive work. Another area of particular interest for future description is to have more detailed descriptions of the few languages where different forms of directional extensions are used for intransitive and transitive predicates (Chapter \ref{ch:identifying2}, Section \ref{sec:allomorphy}). This type of morphology is attested in associated motion paradigms, but not frequently discussed in regard to directional extensions.

The quantitative results presented in Chapter \ref{ch:distribution} give a clear picture of the distribution of directional extensions in Chadic languages which provide a set of general expectations for anyone describing the verbal morphology of Chadic languages. Directional extensions are common, but not ubiquitous. Ventive directional extensions are widespread. Other types of directional extensions are mostly limited to Biu-Mandara. Most languages with directional extensions have only one or two types, usually ventive and itive. In the minority of languages with three or more types of directional extensions (all from Biu-Mandara), less common types, such as bushward and `onto' directions, are attested, while the more common ventive type may actually be absent. 

In the large-scale typological approach of Grambank \citep{thegrambankconsortiumGrambankV12023}, directional verbal morphology is addressed by the single binary question ``Is there directional or locative morphological marking on verbs?'' By comparison, this study of directional verbal morphology in Chadic language is fine-grained and considers variation in the semantic types of directional extensions and their multi-functionality. \citet{martenParametersMorphosyntacticVariation2007} use the term ``micro-variation'' to describe the study of morphosyntactic variation within the context of a family of languages with many shared morphosyntactic features. A statistical approach to micro-variation is what allows this study of directional extensions to affirm the classification of Biu-Mandara languages as a distinct branch. The sub-classification of Chadic languages relies primarily on comparing lexical items. 

Observations about classification naturally raise diachronic questions about how these differences came to be over time. As \citet[5]{schuhChadicCornucopia2017} observes, ``comparative Chadic studies are at an early stage... Aside from some broad areas of agreement, such as the genetic unity of West Chadic or the reconstructability of grammatical gender, there is no `received knowledge' that everyone agrees on.'' In the case of West Chadic languages, the patterns observed suggest that ventive directional extensions are an old type of directional extensions. This analysis is supported by the plausible reconstruction of ventive directional extensions for West Chadic languages (Chapter \ref{ch:diachrony}). However, the fact that ventive directional extensions cannot be reconstructed for Proto-Chadic leaves open the possibility that their frequent appearance could be due to areal influence and might not be a reliable indicator of genetic affiliation. This also applies to the parallels found in other Afroasiatic families, Berber and Semitic, where ventive forms are common \citep{belkadiAssociatedMotionDeictic2015,belkadiAssociatedMotionConstructions2016,fixSemanticsSemiticVentive2020}.

\section{Motion in grammar}

This book makes three main contributions to the study of motion in grammar. First, the comparison of verb roots and verbal morphology in Chapter \ref{ch:verbs} presents a new perspective on the verb-satellite distinction proposed by \citet{talmyLexicalizationPatternsSemantic1985}. It is demonstrated that within the verbal complex of Chadic languages a complementary distribution tendency can be identified such that certain directional meanings (boundary-crossing and vertical) are either found lexicalized in a verb root or grammaticalized in a verbal extension, but not both. This result provides insights into ways in which grammars can evolve to be more efficient and avoid redundancy. However, the same pattern does not hold for ventive and itive directional extensions. This suggests a fundamental distinction between ventive and itive and other types of directional meaning. This type of distinction would not be apparent in research that conflates all types of path of motion into one category. 

The second contribution to the study of motion in grammar is the analysis in Chapter \ref{ch:interpretation} of how co-expressive directional extensions are interpreted. Directional extensions in Chadic languages are multi-functional (Chapter \ref{ch:functions}). The most common function co-expressed with direction is subsequent associated motion. There are also cases of grammaticalized functions co-expressed with direction, and there are many lexicalized uses of combinations of a verb root and a directional extension. Generally speaking, directional extensions have a directional interpretation when used with predicates that describe motion events, with only a limited number of exceptions. Most of these exceptions are cases where a verbal extensions that expresses direction also co-expresses a grammaticalized non-motion meaning, such as aspect. When directional extensions are used with non-motion verbs, the patterns are much more complex -- even unpredictable. Previous studies have focused on the prevalence of associated motion\index{associated motion} interpretations in this context \citep{belkadiAssociatedMotionDeictic2015,belkadiAssociatedMotionConstructions2016,genettiDirectionAssociatedMotion2020,rossCrosslinguisticSurveyAssociated2021}. However, in Chadic languages, other interpretations are possible, especially aspectual meaning and lexicalized combinations.

Where more than one interpretation is possible, the lexical semantics of the verb root is not always sufficient for predicting the interpretation of the directional extension. This is due (in part) to the flexibility that speakers have to frame events with movement (e.g., shooting, closing) as translocative motion events with a moving figure that is not necessarily an argument of the verb. A straightforward analytical approach to composing the meanings of verbs and directional extensions cannot account for this feature of directional extensions without taking into account pragmatic factors \citep{nikitinaPragmaticFactorsVariation2008}. Across linguistic phyla the picture becomes even more complex. Whereas directional extensions in Chadic languages tend to co-express subsequent associated motion (except for those cases where a ventive directional expresses itive prior motion), directional morphology in other languages may tend to co-express prior motion, concurrent motion or some combination \citep{dryerAssociatedMotionDirectionals2021}. Other languages, like English, do not co-express associated motion with directional markers at all. The reasons for these various patterns are not immediately apparent.  

The third contribution to the study of motion in grammar is the correction of previous suggestions that verbs of acquisition (e.g., `take') should be analyzed as motion verbs (Chapter \ref{ch:CAM}) and that buying and selling can be analyzed as motion events (Chapter \ref{ch:buysell}). Since directional morphemes are used in these contexts in Chadic languages, it is not unreasonable to hypothesize that these verbs describe motion events. However, by looking at how these functions are expressed across Chadic languages, it can be seen that this is not a consistent cross-linguistic feature of Chadic languages. Rather, it is one way that these meanings have been lexicalized. Another common form used to lexicalize these meanings is a valency-increasing form, such as a causative. Likewise, treatments that suggest that the use of causative (or valency-increasing) forms in this context is an example of semantic causation are an over-analysis of the grammar. 

The final point to be made is that this book is designed with the assumption that it likely contains some omissions or errors. In case the need for any substantial changes to the data is discovered, the database has been archived in an online deposit in Zenodo \citep{lovestrandDirectionalsChadic2023} so that a new, updated version can be created. Since the same deposit includes the Python scripts used for the calculations, these can be redone in a straightforward manner. Likewise, the publisher, Language Science Press, was chosen for this book because of their online, open access approach to publishing which allows authors the opportunity to update their publications in new editions. Another reason for choosing an online publisher is that this permits the inclusion of an extended set of appendices. The appendices include an account of how the data for each language included in the study were interpreted, and to allow anyone with expertise in a particular Chadic language to offer constructive feedback to the author on any points that need improvement. 
